\documentclass[12pt,a4paper]{article}
\usepackage[utf8]{inputenc}
\usepackage[russian]{babel}
\usepackage[T2A]{fontenc}
\usepackage{geometry}
\usepackage{listings}
\usepackage{xcolor}
\usepackage{graphicx}
\usepackage{hyperref}
\usepackage{indentfirst}

\geometry{left=2.5cm, right=1.5cm, top=2cm, bottom=2cm}

\lstset{
    basicstyle=\ttfamily\footnotesize,
    breaklines=true,
    breakatwhitespace=false,
    frame=single,
    numbers=left,
    numberstyle=\tiny,
    keywordstyle=\color{blue},
    commentstyle=\color{gray},
    stringstyle=\color{red},
    showstringspaces=false,
    tabsize=4,
    inputencoding=utf8,
    extendedchars=true,
    literate={а}{{\cyra}}1 {б}{{\cyrb}}1 {в}{{\cyrv}}1 {г}{{\cyrg}}1
             {д}{{\cyrd}}1 {е}{{\cyre}}1 {ё}{{\cyryo}}1 {ж}{{\cyrzh}}1
             {з}{{\cyrz}}1 {и}{{\cyri}}1 {й}{{\cyrishrt}}1 {к}{{\cyrk}}1
             {л}{{\cyrl}}1 {м}{{\cyrm}}1 {н}{{\cyrn}}1 {о}{{\cyro}}1
             {п}{{\cyrp}}1 {р}{{\cyrr}}1 {с}{{\cyrs}}1 {т}{{\cyrt}}1
             {у}{{\cyru}}1 {ф}{{\cyrf}}1 {х}{{\cyrh}}1 {ц}{{\cyrc}}1
             {ч}{{\cyrch}}1 {ш}{{\cyrsh}}1 {щ}{{\cyrshch}}1 {ъ}{{\cyrhrdsn}}1
             {ы}{{\cyrery}}1 {ь}{{\cyrsftsn}}1 {э}{{\cyrerev}}1 {ю}{{\cyryu}}1
             {я}{{\cyrya}}1 {А}{{\CYRA}}1 {Б}{{\CYRB}}1 {В}{{\CYRV}}1
             {Г}{{\CYRG}}1 {Д}{{\CYRD}}1 {Е}{{\CYRE}}1 {Ё}{{\CYRYO}}1
             {Ж}{{\CYRZH}}1 {З}{{\CYRZ}}1 {И}{{\CYRI}}1 {Й}{{\CYRISHRT}}1
             {К}{{\CYRK}}1 {Л}{{\CYRL}}1 {М}{{\CYRM}}1 {Н}{{\CYRN}}1
             {О}{{\CYRO}}1 {П}{{\CYRP}}1 {Р}{{\CYRR}}1 {С}{{\CYRS}}1
             {Т}{{\CYRT}}1 {У}{{\CYRU}}1 {Ф}{{\CYRF}}1 {Х}{{\CYRH}}1
             {Ц}{{\CYRC}}1 {Ч}{{\CYRCH}}1 {Ш}{{\CYRSH}}1 {Щ}{{\CYRSHCH}}1
             {Ъ}{{\CYRHRDSN}}1 {Ы}{{\CYRERY}}1 {Ь}{{\CYRSFTSN}}1
             {Э}{{\CYREREV}}1 {Ю}{{\CYRYU}}1 {Я}{{\CYRYA}}1
             {→}{{$\rightarrow$}}1
}

\begin{document}

\begin{titlepage}
\begin{center}
    \large
    Национальный исследовательский университет ИТМО\\
    Факультет Программной Инженерии и Компьютерной техники\\
    \vspace{3cm}
    \textbf{\Large ОПИ}\\
    \vspace{0.5cm}
    \textbf{\Large Отчёт по лабораторной работе №2}\\
    \vspace{0.5cm}
    \textbf{\Large Вариант 771267}\\
    \vspace{3cm}
    \begin{flushright}
        Выполнили:\\
        Калинин Дмитрий Викторович\\
        Никифоров Кирилл Евгеньевич\\
        \vspace{0.5cm}
        Группа: P3209\\
        \vspace{0.5cm}
        Преподаватель:\\
        Ермаков Михаил Константинович\\
        Клименков Сергей Викторович\\
    \end{flushright}
    \vfill
    Санкт-Петербург, 2026
\end{center}
\end{titlepage}

\tableofcontents
\newpage

\section{Задание}

\begin{figure}[h!]
    \centering
    \includegraphics[width=\textwidth]{commits.png}
    \caption{Блок-схема варианта лабораторной работы (вариант 771267)}
\end{figure}

Сконфигурировать в своём домашнем каталоге репозитории svn и git и загрузить в них начальную ревизию файлов с исходными кодами (в соответствии с выданным вариантом).

Воспроизвести последовательность команд для систем контроля версий svn и git, осуществляющих операции над исходным кодом, приведённые на блок-схеме.

При составлении последовательности команд необходимо учитывать следующие условия:
\begin{itemize}
    \item Цвет элементов схемы указывает на пользователя, совершившего действие (красный --- первый, синий --- второй).
    \item Цифры над узлами --- номер ревизии. Ревизии создаются последовательно.
    \item Необходимо разрешать конфликты между версиями, если они возникают.
\end{itemize}

\newpage
\section{Создание и конфигурация репозиториев}

\subsection{Git}

Инициализация репозитория:
\begin{lstlisting}[language=bash]
mkdir -p ~/Documents/opi/opi2/git_repo
cd ~/Documents/opi/opi2/git_repo
git init
git config --global user.name "Kirill"
git config --global user.email "kirill@itmo.ru"
\end{lstlisting}

\subsection{SVN}

Создание репозитория и базовой структуры:
\begin{lstlisting}[language=bash]
svnadmin create ~/Documents/opi/opi2/svn_repo
svn mkdir file:///home/ogultra/Documents/opi/opi2/svn_repo/trunk \
          file:///home/ogultra/Documents/opi/opi2/svn_repo/branches \
          --username Kirill --no-auth-cache \
          -m "Init repository structure"
svn checkout file:///home/ogultra/Documents/opi/opi2/svn_repo/trunk \
             ~/Documents/opi/opi2/svn_wc \
             --username Kirill --no-auth-cache
\end{lstlisting}

\newpage
\section{Скрипт для Git}

Скрипт воспроизводит полное дерево из 77 ревизий по блок-схеме варианта.
Красный пользователь --- Kirill, синий --- Dmitry.
Конфликты при слиянии разрешаются автоматически стратегией \texttt{--theirs}
(принимается версия из входящей ветки).

\begin{lstlisting}[language=bash]
#!/bin/bash
set -e

BASE_DIR="$HOME/Documents/opi/opi2"
COMMITS_DIR="$BASE_DIR/commits"
GIT_REPO="$BASE_DIR/git_repo"

rm -rf "$GIT_REPO"
mkdir -p "$GIT_REPO"
cd "$GIT_REPO"
git init

# Kirill = красный (первый пользователь)
# Dmitry = синий  (второй пользователь)

commit_git() {
    local rev=$1
    local author=$2
    local email="${author,,}@itmo.ru"
    cp -rf "$COMMITS_DIR/commit$rev/"*.java . 2>/dev/null
    git add .
    GIT_AUTHOR_NAME="$author"  GIT_AUTHOR_EMAIL="$email" \
    GIT_COMMITTER_NAME="$author" GIT_COMMITTER_EMAIL="$email" \
    git commit -m "Revision r$rev" --allow-empty
}

merge_git() {
    local branch=$1
    local msg=$2
    local author=$3
    local email="${author,,}@itmo.ru"
    GIT_AUTHOR_NAME="$author"  GIT_AUTHOR_EMAIL="$email" \
    GIT_COMMITTER_NAME="$author" GIT_COMMITTER_EMAIL="$email" \
    git merge "$branch" --no-ff -m "$msg" 2>/dev/null || {
        git checkout --theirs -- . 2>/dev/null || true
        git add .
        GIT_AUTHOR_NAME="$author"  GIT_AUTHOR_EMAIL="$email" \
        GIT_COMMITTER_NAME="$author" GIT_COMMITTER_EMAIL="$email" \
        git commit -m "$msg (conflict resolved)"
    }
}

# ШАГ 1. master: r0
commit_git 0 "Kirill"
R0=$(git rev-parse HEAD)

# ШАГ 2. branch_blue6: r0->r1->r11->r18..r21->r35->r37->r52
git checkout -b branch_blue6 $R0
commit_git 1 "Dmitry"; R1=$(git rev-parse HEAD)
commit_git 11 "Dmitry"; R11=$(git rev-parse HEAD)
for i in 18 19 20 21 35 37 52; do commit_git $i "Dmitry"; done

# ШАГ 3. branch_blue2: r1->r2->r22->r40->r42->r50->r57
git checkout -b branch_blue2 $R1
commit_git 2 "Dmitry"; R2=$(git rev-parse HEAD)
for i in 22 40 42 50 57; do commit_git $i "Dmitry"; done

# ШАГ 4. branch_red3: r2->r3->r25->r30->r49->r69
git checkout -b branch_red3 $R2
commit_git 3 "Kirill"; R3=$(git rev-parse HEAD)
for i in 25 30 49 69; do commit_git $i "Kirill"; done

# ШАГ 5. branch_blue7: r3->r4->r17->r27
git checkout -b branch_blue7 $R3
commit_git 4 "Dmitry"; R4=$(git rev-parse HEAD)
commit_git 17 "Dmitry"
commit_git 27 "Dmitry"; R27=$(git rev-parse HEAD)

# ШАГ 6. branch_blue1: r27->r28->r31->r43->r47->r51->r62
git checkout -b branch_blue1 $R27
for i in 28 31 43 47 51 62; do commit_git $i "Dmitry"; done

# ШАГ 7. branch_blue4: r11->r12->r13
git checkout -b branch_blue4 $R11
commit_git 12 "Dmitry"
commit_git 13 "Dmitry"; R13=$(git rev-parse HEAD)

# ШАГ 8. branch_blue3: r13->r14
git checkout -b branch_blue3 $R13
commit_git 14 "Dmitry"

# ШАГ 9. branch_blue4 продолжение: r24->r26->r29->r32->r66
git checkout branch_blue4
for i in 24 26 29 32 66; do commit_git $i "Dmitry"; done

# ШАГ 10. branch_red4: r4->r5->r10->r15
git checkout -b branch_red4 $R4
commit_git 5 "Kirill"; R5=$(git rev-parse HEAD)
commit_git 10 "Kirill"
commit_git 15 "Kirill"; R15=$(git rev-parse HEAD)

# ШАГ 11. branch_red5: r5->r6->r53->r55
git checkout -b branch_red5 $R5
for i in 6 53 55; do commit_git $i "Kirill"; done

# ШАГ 12. branch_red2: r15->r16->r33->r39->r45
git checkout -b branch_red2 $R15
for i in 16 33 39 45; do commit_git $i "Kirill"; done

# ШАГ 13. branch_red4: r36->r38->r54, r58=merge r57, r68, r70=merge r69
git checkout branch_red4
for i in 36 38 54; do commit_git $i "Kirill"; done
merge_git branch_blue2 "Revision r58" "Kirill"
commit_git 68 "Kirill"
merge_git branch_red3 "Revision r70" "Kirill"

# ШАГ 14. branch_blue6: r56=merge r55, r60
git checkout branch_blue6
merge_git branch_red5 "Revision r56" "Dmitry"
commit_git 60 "Dmitry"

# ШАГ 15. branch_blue3: r61=merge r60, r63=merge r62, r73
git checkout branch_blue3
merge_git branch_blue6 "Revision r61" "Dmitry"
merge_git branch_blue1 "Revision r63" "Dmitry"
commit_git 73 "Dmitry"

# ШАГ 16. branch_blue7: r46->r48->r64->r65->r72, r74=merge r73
git checkout branch_blue7
for i in 46 48 64 65 72; do commit_git $i "Dmitry"; done
merge_git branch_blue3 "Revision r74" "Dmitry"

# ШАГ 17. master: r7->r9->r44->r59
git checkout master
commit_git 7 "Kirill"; R7=$(git rev-parse HEAD)
for i in 9 44 59; do commit_git $i "Kirill"; done

# ШАГ 18. branch_blue5: r7->r8->r23->r34->r41, r75=merge r74
git checkout -b branch_blue5 $R7
for i in 8 23 34 41; do commit_git $i "Dmitry"; done
merge_git branch_blue7 "Revision r75" "Dmitry"

# ШАГ 19. branch_red2: r67=merge r66, r71=merge r70, r76=merge r75
git checkout branch_red2
merge_git branch_blue4 "Revision r67" "Kirill"
merge_git branch_red4  "Revision r71" "Kirill"
merge_git branch_blue5 "Revision r76" "Kirill"

# ШАГ 20. master: r77=merge branch_red2
git checkout master
merge_git branch_red2 "Revision r77" "Kirill"

echo ">>> Done!"
\end{lstlisting}

\newpage
\section{Скрипт для SVN}

В SVN ветки реализованы как копии директорий (\texttt{svn copy}).
Структура репозитория: \texttt{trunk} (основная ветка) и \texttt{branches/} (все остальные ветки).
Для гарантии непустых коммитов используется файл-маркер \texttt{.revision}.

\begin{lstlisting}[language=bash]
#!/bin/bash
set -e

BASE_DIR="$HOME/Documents/opi/opi2"
COMMITS_DIR="$BASE_DIR/commits"
SVN_REPO="$BASE_DIR/svn_repo"
SVN_URL="file://$SVN_REPO"
WC="$BASE_DIR/svn_wc"

rm -rf "$SVN_REPO" "$WC"
svnadmin create "$SVN_REPO"

# Хук для установки автора без пароля
cat > "$SVN_REPO/hooks/pre-revprop-change" << 'EOF'
#!/bin/sh
exit 0
EOF
chmod +x "$SVN_REPO/hooks/pre-revprop-change"

get_head_rev() { svnlook youngest "$SVN_REPO"; }

commit_svn() {
    local rev=$1 author=$2
    cp -f "$COMMITS_DIR/commit$rev/"*.java "$WC/" 2>/dev/null
    svn add "$WC/"*.java --force -q 2>/dev/null || true
    echo "r$rev" > "$WC/.revision"
    svn add "$WC/.revision" --force -q 2>/dev/null || true
    svn commit "$WC" --username "$author" --no-auth-cache \
        -m "Revision r$rev"
    svnadmin setrevprop "$SVN_REPO" -r "$(get_head_rev)" \
        svn:author "$author" 2>/dev/null || true
}

branch_svn() {
    local src_url=$1 dst_name=$2 author=$3 msg=$4
    svn copy "$src_url" "$SVN_URL/branches/$dst_name" \
        --username "$author" --no-auth-cache -m "$msg"
    svnadmin setrevprop "$SVN_REPO" -r "$(get_head_rev)" \
        svn:author "$author" 2>/dev/null || true
    svn update "$WC" --no-auth-cache -q 2>/dev/null || true
}

switch_svn() {
    local target=$1
    if [ "$target" = "trunk" ]; then
        svn switch "$SVN_URL/trunk" "$WC" --ignore-ancestry -q
    else
        svn switch "$SVN_URL/branches/$target" "$WC" \
            --ignore-ancestry -q
    fi
    svn update "$WC" --no-auth-cache -q 2>/dev/null || true
}

merge_svn() {
    local src_name=$1 rev_label=$2 author=$3
    local src_url="$SVN_URL/branches/$src_name"
    [ "$src_name" = "trunk" ] && src_url="$SVN_URL/trunk"
    svn update "$WC" --no-auth-cache -q 2>/dev/null || true
    svn merge --allow-mixed-revisions "$src_url" "$WC" -q \
        2>/dev/null || true
    svn resolve --accept working -R "$WC" -q 2>/dev/null || true
    echo "merge-$rev_label" > "$WC/.revision"
    svn commit "$WC" --username "$author" --no-auth-cache \
        -m "Revision $rev_label (merge from $src_name)" || true
    svnadmin setrevprop "$SVN_REPO" -r "$(get_head_rev)" \
        svn:author "$author" 2>/dev/null || true
}

# Инициализация
svn mkdir "$SVN_URL/trunk" "$SVN_URL/branches" \
    --username Kirill --no-auth-cache \
    -m "Init repository structure"
svn checkout "$SVN_URL/trunk" "$WC" \
    --username Kirill --no-auth-cache -q

# r0 - trunk (Kirill)
commit_svn 0 "Kirill"; R0_SVN=$(get_head_rev)

# branch_blue6: r1->r11->r18..r21->r35->r37->r52
branch_svn "$SVN_URL/trunk@$R0_SVN" "branch_blue6" \
    "Dmitry" "Create branch_blue6 from trunk (r0)"
switch_svn "branch_blue6"
commit_svn 1 "Dmitry"; R1_SVN=$(get_head_rev)
commit_svn 11 "Dmitry"; R11_SVN=$(get_head_rev)
for i in 18 19 20 21 35 37 52; do commit_svn $i "Dmitry"; done

# branch_blue2: r2->r22->r40->r42->r50->r57
branch_svn "$SVN_URL/branches/branch_blue6@$R1_SVN" \
    "branch_blue2" "Dmitry" "Create branch_blue2 (r1)"
switch_svn "branch_blue2"
commit_svn 2 "Dmitry"; R2_SVN=$(get_head_rev)
for i in 22 40 42 50 57; do commit_svn $i "Dmitry"; done

# branch_red3: r3->r25->r30->r49->r69
branch_svn "$SVN_URL/branches/branch_blue2@$R2_SVN" \
    "branch_red3" "Kirill" "Create branch_red3 (r2)"
switch_svn "branch_red3"
commit_svn 3 "Kirill"; R3_SVN=$(get_head_rev)
for i in 25 30 49 69; do commit_svn $i "Kirill"; done

# branch_blue7: r4->r17->r27
branch_svn "$SVN_URL/branches/branch_red3@$R3_SVN" \
    "branch_blue7" "Dmitry" "Create branch_blue7 (r3)"
switch_svn "branch_blue7"
commit_svn 4 "Dmitry"; R4_SVN=$(get_head_rev)
commit_svn 17 "Dmitry"
commit_svn 27 "Dmitry"; R27_SVN=$(get_head_rev)

# branch_blue1: r28->r31->r43->r47->r51->r62
branch_svn "$SVN_URL/branches/branch_blue7@$R27_SVN" \
    "branch_blue1" "Dmitry" "Create branch_blue1 (r27)"
switch_svn "branch_blue1"
for i in 28 31 43 47 51 62; do commit_svn $i "Dmitry"; done

# branch_blue4: r12->r13
branch_svn "$SVN_URL/branches/branch_blue6@$R11_SVN" \
    "branch_blue4" "Dmitry" "Create branch_blue4 (r11)"
switch_svn "branch_blue4"
commit_svn 12 "Dmitry"
commit_svn 13 "Dmitry"; R13_SVN=$(get_head_rev)

# branch_blue3: r14
branch_svn "$SVN_URL/branches/branch_blue4@$R13_SVN" \
    "branch_blue3" "Dmitry" "Create branch_blue3 (r13)"
switch_svn "branch_blue3"
commit_svn 14 "Dmitry"

# branch_blue4 продолжение: r24->r26->r29->r32->r66
switch_svn "branch_blue4"
for i in 24 26 29 32 66; do commit_svn $i "Dmitry"; done

# branch_red4: r5->r10->r15
branch_svn "$SVN_URL/branches/branch_blue7@$R4_SVN" \
    "branch_red4" "Kirill" "Create branch_red4 (r4)"
switch_svn "branch_red4"
commit_svn 5 "Kirill"; R5_SVN=$(get_head_rev)
commit_svn 10 "Kirill"
commit_svn 15 "Kirill"; R15_SVN=$(get_head_rev)

# branch_red5: r6->r53->r55
branch_svn "$SVN_URL/branches/branch_red4@$R5_SVN" \
    "branch_red5" "Kirill" "Create branch_red5 (r5)"
switch_svn "branch_red5"
for i in 6 53 55; do commit_svn $i "Kirill"; done

# branch_red2: r16->r33->r39->r45
branch_svn "$SVN_URL/branches/branch_red4@$R15_SVN" \
    "branch_red2" "Kirill" "Create branch_red2 (r15)"
switch_svn "branch_red2"
for i in 16 33 39 45; do commit_svn $i "Kirill"; done

# branch_red4: r36->r38->r54, r58=merge r57, r68, r70=merge r69
switch_svn "branch_red4"
for i in 36 38 54; do commit_svn $i "Kirill"; done
merge_svn "branch_blue2" "r58" "Kirill"
commit_svn 68 "Kirill"
merge_svn "branch_red3" "r70" "Kirill"

# branch_blue6: r56=merge r55, r60
switch_svn "branch_blue6"
merge_svn "branch_red5" "r56" "Dmitry"
commit_svn 60 "Dmitry"

# branch_blue3: r61=merge r60, r63=merge r62, r73
switch_svn "branch_blue3"
merge_svn "branch_blue6" "r61" "Dmitry"
merge_svn "branch_blue1" "r63" "Dmitry"
commit_svn 73 "Dmitry"

# branch_blue7: r46->r48->r64->r65->r72, r74=merge r73
switch_svn "branch_blue7"
for i in 46 48 64 65 72; do commit_svn $i "Dmitry"; done
merge_svn "branch_blue3" "r74" "Dmitry"

# trunk: r7, branch_blue5, r9->r44->r59
switch_svn "trunk"
commit_svn 7 "Kirill"; R7_SVN=$(get_head_rev)
branch_svn "$SVN_URL/trunk@$R7_SVN" "branch_blue5" \
    "Dmitry" "Create branch_blue5 from trunk (r7)"
for i in 9 44 59; do commit_svn $i "Kirill"; done

# branch_blue5: r8->r23->r34->r41, r75=merge r74
switch_svn "branch_blue5"
for i in 8 23 34 41; do commit_svn $i "Dmitry"; done
merge_svn "branch_blue7" "r75" "Dmitry"

# branch_red2: r67=merge r66, r71=merge r70, r76=merge r75
switch_svn "branch_red2"
merge_svn "branch_blue4" "r67" "Kirill"
merge_svn "branch_red4"  "r71" "Kirill"
merge_svn "branch_blue5" "r76" "Kirill"

# trunk: r77=merge branch_red2
switch_svn "trunk"
merge_svn "branch_red2" "r77" "Kirill"

echo ">>> SVN Done!"
\end{lstlisting}

\newpage
\section{Вывод}

В ходе выполнения лабораторной работы были настроены репозитории SVN и Git
и воспроизведена последовательность из 77 ревизий в соответствии с блок-схемой варианта.

Были изучены основные команды для работы с системами контроля версий:
\texttt{commit}, \texttt{merge}, \texttt{checkout}, \texttt{branch} для Git;
\texttt{commit}, \texttt{copy}, \texttt{switch}, \texttt{merge}, \texttt{resolve} для SVN.

Изучены механизмы ветвления и способы разрешения конфликтов при слиянии веток,
возникающих при одновременном редактировании файлов разными пользователями.
Конфликты разрешались как вручную через редактор vim (удаление маркеров \texttt{<<<<<<<},
\texttt{=======}, \texttt{>>>>>>>}), так и автоматически через стратегию \texttt{--theirs}.

Полученные навыки применимы при командной разработке программного обеспечения.

\end{document}
